\scanpage{142}
% archive foliation: 141

\uline{Le foncteur \(\unsure{\varphi|\Phi}\) est pleinement fidèle}
(On se ramène au cas de la base \(\eta=\operatorname{Spec}K\)).

L'image essentielle est \deleted{\uline{\unsure{évidemment|exactement}} \unsure{donc}}
formée des motifs \((M,\{E_l\},\{u_l\})\) tels que
\deleted{les \(E_l\)} leur quotient modulo
torsion est associé à un \uline{schéma} \uline{abélien} sur \(S\).
Il \uline{faut} pour cela que \(M\) soit effectif de poids \(1\),
i.e. est défini par un schéma abélien \inserted{\(A_K\)} sur \(K\).
\deleted{\(\unsure{A|R}\)} \deleted{\(M\)}.
La question si cette condition est suffisante revient à
celle-ci: Soit \(A\) un schéma abélien sur \(K\), \deleted{\unsure{de}} tel
que \(\forall l\), \deleted{les} \(T_l(A)\) soit non ramifié sur \(S_l\), alors
\(A\) a bonne réduction sur \(S\). Cela \unsure{marche} si
\(S\) est de dim \(1\), \inserted{(critère de \inserted{Néron-}Ogg-\unsure{Chafarevitch})},
\uline{ou} si de car. nulle (critère de
\unsure{critique|néronique}), mais peut-être pas au général [précisément
local régulier de dim \(2\) d'inégales caractéristiques]. Pour
faire un contre-exemple, il suffirait d'avoir un exemple
d'un schéma abélien \(A\) sur \(S\), et d'un sous-groupe fini de
\(A_K\) qui ne se prolonge pas en un sous-groupe fini et \uline{plat}
de \(A\). S'il y avait un tel exemple, il y en aurait
aussi avec un indice premier......

\margin{\uline{Pb} Avec la définition adoptée de motif lisse (en termes de
celle de isomotif lisse, comme un isomotif \uline{sur} \(K\)
avec simplement une propriété \inserted{``de lissité''} de non ramification pour
les \(T_l\)) il n'est pas du tout clair comment définir
l'image inverse d'un motif (resp. d'un isomotif) par
\(S'\mathbin{-}S\), p.ex. les fibres en les \(s\in S\). Cette difficulté
\deleted{se pose}\inserted{concerne}
--- tous les isomotifs, et vient déjà pour ceux qui sont effectifs de
poids \(1\).}
